\begin{table}[tbp]
\caption{References, validation, and permissible claims.}
\label{tab:ch3-validation}
\begingroup\fontsize{9.2}{11.8}\selectfont
\renewcommand{\arraystretch}{1.12}
\begin{tabularx}{\textwidth}{@{}>{\raggedright\arraybackslash}p{25mm}>{\raggedright\arraybackslash}X >{\raggedright\arraybackslash}X@{}}
\toprule
\textbf{Comparison} & \textbf{Reference and supporting evidence} & \textbf{Scope of the conclusion}\\
\midrule
Implementation agreement & Identical inputs and parameters; an independently checked computation or retained reference output. & Agreement in realizing the specified measure or model. Semantic validity requires additional evidence (Chs.~4--6).\\[5pt]
Score comparison & Other measures and participant assessments, with the compared population and shared inputs identified. & Correspondence or disagreement within that population; not an independent contribution-quality label (Ch.~5).\\[5pt]
Simulated diffusion & Common graph, seed budget and propagation model; reported reach for ranked seed sets. & Performance in the defined simulation; real diffusion and benefit require further observation (Ch.~5).\\[5pt]
Prediction & Outcomes unavailable during fitting and selection; splits appropriate to time and transfer claims. & Performance on the stated target and test population. Predicting an index validates neither its meaning nor its use.\\[5pt]
Field deployment & Dated introduction, pre-launch history, a justified reference, participant accounts and observation checks. & Evidence about the deployed arrangement under recorded conditions; stronger causal claims depend on identification (Ch.~6).\\
\bottomrule
\end{tabularx}
\endgroup
\tablenote{Methodological synthesis. The completed studies supply different subsets of this evidence; the table does not attribute every recommended check to their original protocols.}
\end{table}
